\begin{statement}{E1}{Proved}{Vertical divisors and the field of meromorphic functions}
Let $Y$ be a connected compact complex manifold of dimension $n\geq2$
with $H^2(Y,\mathbb R)=0$. Let $g:Y\to C$ be a surjective holomorphic
map to a connected compact Riemann surface, with connected fibres.
Suppose that its smooth fibres over a nonempty open subset of $C$ are
K\"ahler. Every effective divisor of $Y$ is supported on finitely many
fibres, and pullback induces an isomorphism of fields
\[
 g^*: \mathcal M(C)\xrightarrow{\ \sim\ }\mathcal M(Y).
\]
In particular $a(Y)=1$. Reducedness of the other fibres is unnecessary.
\end{statement}
\begin{proof}
For a prime divisor $D$, the image $g(D)$ is analytic by properness. It
is a point or all of $C$. In the latter case choose a regular value $c$
in the stated open subset for which $F=g^{-1}(c)$ is not contained in
$D$. The canonical section of $\mathcal O_Y(D)$ restricts to a nonzero
section with nonempty effective zero divisor $D_F$ on $F$. For a
K\"ahler form $\omega$ on $F$, the divisor class formula gives
\[
 0=\int_F c_1(\mathcal O_Y(D)|_F)\wedge\omega^{n-2}
   =\int_{D_F}\omega^{n-2}>0.
\]
The first equality follows from the restriction homomorphism
$H^2(Y,\mathbb R)\to H^2(F,\mathbb R)$; the last follows from positivity
on every component with its positive multiplicity. Thus $g(D)$ is a
point. A divisor has finitely many prime components, proving the first
assertion.

For $h\in\mathcal M(Y)$ remove from $C$ the finitely many values
supporting its polar divisor. On the inverse image of the complement,
$h$ is holomorphic: a meromorphic function on a manifold having no
codimension-one pole extends across its codimension-two indeterminacy
set by the local normality of the structure sheaf. Its restriction to
each compact connected fibre is constant. The identity
$g_*\mathcal O_Y=\mathcal O_C$, supplied by Stein factorization and
the normality of $C$, therefore gives a holomorphic descended function
there. To check meromorphic extension at the removed values, take the
closure of the graph of $h$ in $Y\times\mathbb P^1$ and apply the proper
map $g\times1$. Its image is a one-dimensional analytic subset of
$C\times\mathbb P^1$, generically the graph of the descended function.
Its component dominating $C$ has generic degree one, and hence is the
graph of a meromorphic function on $C$. Pullback equals $h$ on a dense
open subset and therefore in the meromorphic function field. Pullback
is injective because $g$ is surjective. This proves the field
isomorphism and its asserted transcendence degree.
\end{proof}


\subsection{E.2. Nonholomorphic reduction and the printed proof}

% E2 from part_E.tex

